\documentclass[12pt,a4paper]{article}
\usepackage[utf8]{inputenc}
\usepackage[T1]{fontenc}
\usepackage{lmodern}
\usepackage[margin=2.5cm]{geometry}
\usepackage{setspace}
\onehalfspacing

\begin{document}

\noindent
2 April 2026

\bigskip

\noindent
Dear Editor-in-Chief,\\
\textit{Stochastic Environmental Research and Risk Assessment}

\bigskip

\noindent
We are pleased to submit the manuscript entitled ``Analytical uncertainty of the statistical paradigm as a stochastic risk vector in soil quality assessment under long-term conservation management'' for consideration as an original research article in \textit{Stochastic Environmental Research and Risk Assessment}.

\medskip

\noindent
Soil quality indices are widely used to guide conservation recommendations, yet each study selects a single scoring algorithm and implicitly assumes that the resulting ranking is method-independent. We challenge this assumption by benchmarking five structurally distinct paradigms (ANOVA/GLM, fuzzy inference, PLS-SEM, MARS, and Random Forest) on the same 22-year field experiment ($n = 108$, 18 indicators, four edaphic domains, Typic Hapludult, northeast Brazil). The main contributions include: (i) formal concordance assessment across paradigms using Kendall $W$, ICC, and Bland--Altman limits of agreement; (ii) demonstration that nonlinear methods form a near-interchangeable block ($\rho > 0.99$, LoA $< \pm 0.43$), whereas linear paradigms diverge systematically; (iii) Monte Carlo evidence that inter-paradigm concordance stabilizes from $n \geq 30$ when chemical indicators are included; and (iv) a decision flowchart linking analytical priority to paradigm choice with explicit uncertainty thresholds.

\medskip

\noindent
To the best of our knowledge, no previous study has simultaneously compared five paradigms from different analytical families on the same edaphic dataset using formal stochastic concordance metrics. The manuscript therefore addresses a relevant gap at the interface of pedometrics and environmental risk assessment.

\medskip

\noindent
We confirm that this work is original and is not under consideration elsewhere. All authors have read and approved the final version. We declare no conflicts of interest. Thank you for the opportunity to present our work.

\bigskip

\noindent
Sincerely,

\bigskip\bigskip

\noindent
Luiz Diego Vidal Santos\\
Universidade Estadual de Feira de Santana (UEFS), Brazil\\
ldvsantos@uefs.br\\
ORCID: 0000-0001-8659-8557

\end{document}
